\begin{frame}{Current Julia AC-OPF Model: Formulation}
\small
\begin{itemize}
  \item Bus-injection formulation: \code{Pg}, \code{Qg}, \code{Vm}, \code{Va}.
  \item Power balance uses the sparse row of \code{Ybus}.
\end{itemize}

\[
P_i =
\sum_j V_iV_j\left(G_{ij}\cos(\theta_i-\theta_j)+B_{ij}\sin(\theta_i-\theta_j)\right)
\]
\[
Q_i =
\sum_j V_iV_j\left(G_{ij}\sin(\theta_i-\theta_j)-B_{ij}\cos(\theta_i-\theta_j)\right)
\]

\begin{itemize}
  \item Bus shunts are already in the diagonal:
\end{itemize}
\[
Y_{ii} = \sum_{\ell \in \delta(i)} Y_{ii}^{(\ell)} + g_i^{sh} + j b_i^{sh}
\qquad
Y_{ij} = \sum_{\ell:\, i \leftrightarrow j} Y_{ij}^{(\ell)}
\]

\begin{itemize}
  \item Shunt susceptance appears through the \(-B_{ii}V_i^2\) part of \code{Qi}.
  \item Parallel lines accumulate in the same off-diagonal \code{Yij}.
  \item Thermal/rate limits still stay per physical branch.
\end{itemize}

\vspace{0.2em}
\tbl{\linkFormulationSupport}{support folder}
\tbl{\linkBusShuntAudit}{bus shunts}
\tbl{\linkParallelBranchAudit}{parallel lines}
\tbl{\linkShuntToggleSummary}{shunt toggle}

\end{frame}


\begin{frame}{AC-OPF Timing: Julia and Python PM-like}
\small
\begin{itemize}
  \item Current: bus-injection expressions directly in balances and thermal limits.
  \item PM-like: explicit branch-flow variables plus Ohm's-law flow constraints.
  \item Cell format: mean solve time / mean iterations.
\end{itemize}

\vspace{0.2em}
\begin{tabular}{lllll}
\toprule
Case & Julia current & Julia PM-like & Python PM-like & PowerModels \\
\midrule
5\_pjm & 0.161 / 19.5 & 0.029 / 21.1 & 0.096 / 21.3 & 0.016 / 21.1 \\
14\_ieee & 0.014 / 15.1 & 0.012 / 15.4 & 0.099 / 15.5 & 0.009 / 15.4 \\
30\_ieee & 0.024 / 16.8 & 0.022 / 17.5 & 0.112 / 17.6 & 0.020 / 17.5 \\
118\_ieee & 0.115 / 25.2 & 0.111 / 28.3 & 0.274 / 28.3 & 0.099 / 28.3 \\
300\_ieee & 0.348 / 32.2 & 0.342 / 34.1 & 0.617 / 34.1 & 0.310 / 34.1 \\
1354\_pegase & 2.909 / 40.9 & 2.841 / 50.1 & 3.974 / 50.3 & 2.812 / 50.1 \\
1888\_rte & 15.437 / 204.3 & 10.276 / 131.0 & 13.091 / 143.1 & 11.293 / 155.9 \\
6470\_rte & 128.685 / 427.6 & 27.755 / 94.4 & 30.322 / 94.4 & 27.687 / 94.2 \\
\bottomrule
\end{tabular}

\vspace{0.3em}
\begin{itemize}
  \item Python PM-like iteration counts added from the Ipopt iteration backfill.
\end{itemize}

\tbl{\linkPerfAllModelsPythonPMlike}{all-model timing}
\tbl{\linkPMvsCurrentPerf}{compact timing}
\tbl{\linkPMvsCurrentPerfByCase}{per-case timing}
\end{frame}


\begin{frame}{Deviation Minimization: Summary}
\scriptsize
\begin{itemize}
  \item Scenario filter: non-zero objective gap from existing objective-mismatch diagnostics.
  \item Diagnostic objective: minimize deviation from the PowerModels OPF point.
\end{itemize}

\vspace{0.2em}
\begin{tabular}{lllrrr}
\toprule
Model & Scenario & Cheaper & Obj gap & Full L1 & Full max var \\
\midrule
Current Julia & load x1.2 reactive & current & -3.84e3 & 0 & 6.78e-10 \\
Current Julia & load x0.8 & current & -12.5 & 0 & 5.79e-10 \\
PM-like Julia & load x0.8 & PM & 7.98e4 & 0 & 5.83e-10 \\
PM-like Julia & load x1.2 reactive & PM & 3.58e3 & 0 & 5.97e-10 \\
PM-like Julia & Qg limits x0.6 & PM & 3.93e2 & 0 & 1.90e-10 \\
PM-like Python & load x0.8 & PM & 7.98e4 & 3.86e-4 & 2.88e-5 \\
PM-like Python & load x1.2 reactive & PM & 3.58e3 & 2.22e-4 & 4.88e-5 \\
\bottomrule
\end{tabular}

\vspace{0.4em}
\begin{itemize}
  \item Julia diagnostics reach the PowerModels point to numerical precision.
  \item Python PM-like lands close to PowerModels; residual differences remain small but visible.
\end{itemize}

\tbl{\linkDevMinCurrentSummary}{current summary}
\tbl{\linkDevMinPMlikeSummary}{PM-like summary}
\tbl{\linkDevMinPMlikePythonSummary}{Python summary}
\tbl{\linkDevMinVariantSummary}{variants}
\tbl{\linkObjectiveGapFilteredScenarios}{objective filter}
\tbl{\linkConstrainedDeviationConstraintMismatchByScenario}{constraint scenario}
\tbl{\linkConstrainedDeviationConstraintMismatchByFamily}{constraint family}
\tbl{\linkConstrainedDeviationConstraintMismatchDetails}{constraint details}
\tbl{\linkConstrainedDevNonzeroLOneMismatchByFamily}{nonzero-L1 family}
\tbl{\linkConstrainedDevNonzeroLOneMismatchDetails}{nonzero-L1 details}
\tbl{\linkConstrainedDevNonzeroLOneMaterialMismatchDetails}{material nonzero-L1}
\end{frame}

\begin{frame}{Equality-Only Deviation Problem}
\scriptsize
\begin{itemize}
  \item Keep only AC equality constraints.
  \item Drop inequality constraints during solve; evaluate \(h(\cdot)\) afterwards.
  \item Compare \(h(x_E^{**},u_E^{**})\) against \(h(x_P^*,u_P^*)\).
  \item All-model h-level table includes current Julia, PM-like Julia, and PM-like Python.
\end{itemize}

\vspace{0.2em}
\begin{tabular}{llrrrrl}
\toprule
Model & Scenario & L1 & Max var & h mism. & h viol. & Detail \\
\midrule
current & load x1.2 reactive & 0 & 7.14e-13 & 0 & 0 & \tbl{\linkEqOnlyReactive}{open} \\
current & load x0.8 & 0 & 4.71e-13 & 0 & 0 & \tbl{\linkEqOnlyLoadLow}{open} \\
PM-like & load x0.8 & 0 & 2.28e-13 & 0 & 0 & \tbl{\linkEqOnlyLoadLow}{open} \\
PM-like & load x1.2 reactive & 0 & 4.60e-13 & 0 & 0 & \tbl{\linkEqOnlyReactive}{open} \\
PM-like & Qg limits x0.6 & 0 & 4.32e-13 & 0 & 0 & \tbl{\linkEqOnlyQgLimits}{open} \\
PM-like Python & load x0.8 & 0 & 8.09e-13 & 0 & 0 & \tbl{\linkConstraintMismatchEvidence}{tables} \\
PM-like Python & load x1.2 reactive & 0 & 1.25e-12 & 0 & 0 & \tbl{\linkConstraintMismatchEvidence}{tables} \\
\bottomrule
\end{tabular}

\vspace{0.4em}
\begin{itemize}
  \item Max \(h(\cdot)\) margin differences are \(10^{-12}\)-level in the summary.
  \item Therefore, \(h(x_E^{**},u_E^{**}) = h(x_P^*,u_P^*)\) to numerical precision.
\end{itemize}

\tbl{\linkEqOnlySummary}{summary}
\tbl{\linkEqOnlyHLevelsAllModels}{h levels}
\tbl{\linkEqOnlyReconciliation}{comparison}
\tbl{\linkEqualityOnlyConstraintLevelByScenario}{constraint scenario}
\tbl{\linkEqualityOnlyConstraintLevelByFamily}{constraint family}
\tbl{\linkEqualityOnlyConstraintLevelDetails}{constraint details}
\end{frame}


\begin{frame}{Load-Flow Validation Workflow}
\small
\begin{itemize}
  \item Start from both OPF endpoints: own solution and PowerModels solution.
  \item Warm-start data copied from each OPF endpoint: bus \code{Vm}/\code{Va}, generator \code{Pg}/\code{Qg}, and generator \code{Vg}.
  \item OPF bus voltages provide the Newton initial point; generator dispatch and voltage setpoints define the PF case.
  \item PF convergence then gives the post-PF operating point used for margin checks.
  \item Julia path: PowerModels \code{compute\_ac\_pf}; branch flows recomputed with the same AC branch-flow function.
  \item Python path: pypowsybl/OpenLoadFlow; branch flows taken from the converged pypowsybl PF result.
  \item Compare post-PF inequality margins and binding changes at both endpoints.
  \item Focus table: Julia has post-PF binding-level mismatches but zero endpoint/shared limit violations.
  \item Python PM-like focus rows still show endpoint/shared limit violations for follow-up.
\end{itemize}

\vspace{0.4em}
\tbl{\linkLoadflowWorkflow}{LF validation folder}
\tbl{\linkLoadflowScope}{LF scope}
\tbl{\linkLoadflowEndpoint}{endpoint PF rows}
\tbl{\linkLoadflowMismatchViolationSummary}{focused summary}
\tbl{\linkLoadflowMismatchViolationFocus}{focused details}
\tbl{\linkPythonPMlikeLFViolationDiagnosis}{Python diagnosis}
\tbl{\linkLoadflowConstraintLevelByScenario}{constraint scenario}
\tbl{\linkLoadflowConstraintLevelByFamily}{constraint family}
\tbl{\linkLoadflowConstraintLevelDetails}{constraint details}
\end{frame}
